\chapter{Relations and compatible homotopies}
\label{chap:relations-homotopies}

\section{A differential for a scalar relation}

An abelian group is a set with an associative,
commutative addition, an additive identity $0$, and an additive inverse
$-a$ for every element $a$. A ring is an abelian group with an associative
multiplication and a multiplicative identity $1$, satisfying
$a(b+c)=ab+ac$ and $(a+b)c=ac+bc$.
A ring is commutative when $ab=ba$ for every $a,b$.
All rings and ring homomorphisms here have identities, and the
homomorphisms preserve addition, multiplication, and $1$.

Let $k$ be a field: a commutative ring with $1\ne0$ in which every
nonzero element has a multiplicative inverse. Assume that $k$ has
characteristic zero, meaning that $n\cdot1\ne0$ for every positive
integer $n$. In particular, $2$ is invertible in $k$.

Let $R=k[t]$ consist of finite sums $\sum_{j\geq0}a_jt^j$ with $a_j\in k$.
Two such sums agree when every coefficient agrees.
Addition adds coefficients, and multiplication distributes the rule
$t^it^j=t^{i+j}$ over finite sums.
These operations give a commutative ring, called the polynomial ring.
The symbol $t$ is an indeterminate because it satisfies no further relation.
Fix an integer $r\geq1$ and consider the relation $t^r=0$.

Put $s=t^r$. The quotient $R/(s)$ identifies polynomials $f,g$ when
$f-g$ is a multiple of $s$. Write $[f]$ for the class of $f$.
The operations are $[f]+[g]=[f+g]$ and $[f][g]=[fg]$.
Replacing $f,g$ by $f+sa,g+sb$ changes their sum by $s(a+b)$
and their product by $s(ag+fb+sab)$.
Thus these operations do not depend on the representatives.
Multiplication by $s$ is the zero map on this quotient.


The product of two nonzero polynomials is nonzero: its highest
coefficient is the product of their highest nonzero coefficients in $k$.
The same argument applies after adjoining another polynomial variable.
Thus $k[t,u]=k[t][u]$ also has no zero divisors, meaning that a product
of two nonzero elements cannot be zero.

An $R$-module is an abelian group with a scalar action of $R$.
The action satisfies $a(m+n)=am+an$, $(a+b)m=am+bm$,
$(ab)m=a(bm)$, and $1m=m$.
A free module with basis $v_i$ consists of the finite sums $\sum_i a_i v_i$,
with uniquely determined coefficients $a_i\in R$.
A map of $R$-modules is $R$-linear if it preserves sums and scalar
multiplication.
Its kernel is the set of inputs sent to zero, and its image is the set of
all outputs. Both are submodules: subsets containing zero and closed under
addition, additive inverses, and scalar multiplication.

An associative unital $R$-algebra is an $R$-module with an associative
product, an identity element, and an $R$-bilinear multiplication:
the product distributes over sums and satisfies $(ra)b=r(ab)=a(rb)$.
The scalar map $R\to A$ sends $r$ to $r1$.
It preserves ring operations, and its image is central because
$(r1)a=ra=a(r1)$.

A graded $R$-module is a family of $R$-modules $A^n$, indexed by
integers, with underlying module $A=\bigoplus_n A^n$.
The direct sum consists of elements with finitely many nonzero components.
An element of $A^n$ is homogeneous of degree $n$, written $|a|=n$.

A graded $R$-algebra is such an algebra
$A=\bigoplus_{n\in\mathbb Z}A^n$ whose product satisfies
$A^p A^q\subset A^{p+q}$. Scalars lie in degree zero and commute with
every element. The direct sum requires each element to have finitely many
nonzero homogeneous components.

A linear map $F:V\to W$ between graded modules has degree $n$ if
$F(V^j)\subset W^{j+n}$ for every integer $j$.
\begin{definition}\label{def:rel-dga}
A differential graded $R$-algebra is a graded $R$-algebra with an
$R$-linear differential $d$ of degree one such that $d^2=0$, $d1=0$, and
\begin{equation}\label{eq:rel-leibniz}
 d(ab)=(da)b+(-1)^{|a|}a\,db
\end{equation}
for homogeneous elements $a,b$. This identity is the graded Leibniz rule.
\end{definition}


The sign depends only on the parity of $|a|$, namely its remainder modulo
two. In particular, the sign is negative when $|a|=-1$.
The rule makes the two middle terms in $d^2(ab)$ cancel.
An element $h$ of degree $n-1$ satisfying $dh=v$ is a primitive of
an element $v$ of degree $n$.
In a differential graded $R$-algebra $A$, multiplication by $s$ is
the degree-zero map $x\mapsto sx$.
A null-homotopy for this map is an $R$-linear map $H:A\to A$ of
degree minus one satisfying
\[
 dH(x)+H(dx)=sx\qquad(x\in A).
\]
This asks for an equation between maps on $A$, while the quotient
$R/(s)$ makes scalar multiplication by $s$ identically zero.
If $dx=0$, this equation gives the primitive $H(x)$ of $sx$.
A primitive $\tau\in A^{-1}$ of $s1$ defines left multiplication
$L_\tau(x)=\tau x$, which supplies such a map:
the graded Leibniz rule gives
\[
 d(\tau x)+\tau\,dx=sx.
\]
Thus a primitive of the scalar relation supplies a null-homotopy
on the whole algebra.

We ask whether the null-homotopy can also respect left and right
multiplication. A homogeneous map $h$ of degree $-1$ compatible with left and
right multiplication must satisfy
\begin{equation}\label{eq:rel-homotopy-linearity}
 h(amb)=(-1)^{|a|}a\,h(m)b
\end{equation}
for homogeneous $a,m,b$. The right action contributes no sign in this
convention. The left sign accounts for moving a degree-minus-one map past
the left input. We will derive this convention from a single algebra
action below.

For $h=L_\tau$, right multiplication satisfies this equation because
$\tau(mb)=(\tau m)b$.
The left multiplication condition at $a=\tau$, $m=b=1$ requires
$\tau^2=-\tau^2$.
Hence $L_\tau$ can respect both actions only if $\tau^2=0$,
since $2$ is invertible.
The equation $d\tau=s$ alone does not impose this additional relation.
We must determine the square of a primitive before claiming that its
null-homotopy respects multiplication.

\section{Freely adjoining a primitive}
\label{sec:rel-algebra}

To retain every product permitted by $d\tau=s$, impose no additional
equation on $\tau$.
Let $E=R\langle\tau\rangle$ be the free associative $R$-algebra on one
generator of degree $-1$. Its $R$-basis is
\[
 1,\tau,\tau^2,\tau^3,\ldots,
 \qquad \tau^i\tau^j=\tau^{i+j},\qquad |\tau^i|=-i.
\]
Define
\begin{equation}\label{eq:rel-state-witness}
 dt=0,\qquad d\tau=t^r,\qquad |\tau|=-1,
\end{equation}
and extend by the graded Leibniz rule.

Repeated application of that rule gives
\begin{equation}\label{eq:rel-word-differential}
 d(\tau^n)
 =s\sum_{i=0}^{n-1}(-1)^i\tau^{n-1}
 =\begin{cases}
 0,&n\text{ even},\\
 s\tau^{n-1},&n\text{ odd}.
 \end{cases}
\end{equation}
For $n=0$ the sum is empty and equals zero.
The formula proves $d^2=0$ on each basis element.

Put $u=\tau^2$, so $|u|=-2$ and $du=0$.
The basis shows $u\ne0$, so the map $L_\tau$ fails the left
multiplication condition.
Every element of $E$ has a unique expression $a(u)+\tau b(u)$ with
$a,b\in R[u]$. Equation~\eqref{eq:rel-word-differential} becomes
\[
 E=R[u]\oplus\tau R[u],\qquad
 d\bigl(a(u)+\tau b(u)\bigr)=s b(u).
\]
The two summands consist of even and odd degrees, respectively.

The failure of left compatibility can itself have a primitive after
multiplying the scalar relation by another polynomial.
To compute when this happens, let $\lambda\in R$ and choose a
primitive $h\in E^{-1}$ with $dh=\lambda$.
The map $L_h(x)=hx$ respects right multiplication.
It is a null-homotopy for multiplication by $\lambda$, by the graded
Leibniz rule.
Its left defect on the generator $\tau$ is $\tau h+h\tau$.
We ask for a primitive $b\in E^{-3}$ of that defect:
\[
 dh=\lambda,\qquad db=\tau h+h\tau.
\]
The degree conditions give $h=A\tau$ and $b=B\tau u$ for
polynomials $A,B\in R$.
Since $d(\tau u)=su$, the two equations become
\[
 sA=\lambda,\qquad sB=2A.
\]
They imply $\lambda=s^2B/2$.
Conversely, if $\lambda=s^2c$ with $c\in R$, take
$A=sc$ and $B=2c$.
Thus the two primitive equations have a solution precisely when the
scalar is a multiple of $s^2=t^{2r}$.
This asks the multiplication defect to have a primitive.
Strict compatibility would require the defect itself to be zero.
In particular, $t^2$ has a primitive when $r=2$, but this correction
requires a multiple of $t^4$.

\section{A complex for a polynomial relation}
\label{sec:rel-complexes}

For a submodule $U\subset M$, the quotient $M/U$ identifies $m$ with $m'$
when $m-m'\in U$. Write $[m]$ for the resulting class.
Define $[m]+[n]=[m+n]$ and $a[m]=[am]$.
Replacing $m,n$ by $m+u,n+v$ changes these expressions by $u+v$ and $au$,
which belong to $U$. The operations are therefore well-defined.
They satisfy the module laws because the original operations do.

An additive subgroup is a subset containing zero and closed under
addition and additive inverses.
In a ring $A$, a two-sided ideal is an additive subgroup $I$ such that
$ax,xa\in I$ whenever $a\in A$ and $x\in I$.
The quotient $A/I$ also has multiplication $[a][b]=[ab]$:
replacing $a,b$ by $a+x,b+y$ changes $ab$ by $ay+xb+xy\in I$.
The ideal generated by elements $x_i$ consists of finite sums $a x_i b$.
For a commutative ring these are finite sums $a x_i$.
In particular, $(t^r)\subset R$ consists of the multiples of $t^r$,
and $R/(t^r)$ imposes the relation $t^r=0$.

Division by $t^r$ gives each class a unique representative
\[
 a_0+a_1t+\cdots+a_{r-1}t^{r-1},\qquad a_i\in k.
\]
The quotient records which polynomials agree after imposing the relation.
The two-term sequence
\[
 R\tau\xrightarrow{\ d\ }R,
 \qquad d(a\tau)=at^r,
\]
retains an element $a\tau$ mapped to each multiple $at^r$ of the relation.
We place $R\tau$ in degree $-1$ and $R$ in degree zero.
All other degrees are zero.

\begin{definition}\label{def:rel-complex}
A cochain complex of $R$-modules consists of modules $V^n$, for integers
$n$, and $R$-linear maps $d:V^n\to V^{n+1}$ with $d^2=0$.
Write $V=\bigoplus_{n\in\mathbb Z}V^n$: each element is a sum of
finitely many components in the indicated modules.
An element of $V^n$ is homogeneous of degree $n$, written $|v|=n$.
The cycles, boundaries, and cohomology in degree $n$ are
\[
 \begin{split}
 Z^n(V)&=\ker(d:V^n\longrightarrow V^{n+1}),\\
 B^n(V)&=\operatorname{im}(d:V^{n-1}\longrightarrow V^n),\\
 H^n(V)&=Z^n(V)/B^n(V).
 \end{split}
\]

We write $[v]$ for the cohomology class of a cycle $v$.
The notation $H^*(V)$ means the graded module $\bigoplus_n H^n(V)$.
A complex is acyclic if all its cohomology groups are zero.
An element $h\in V^{n-1}$ with $dh=v$ is a primitive of $v\in V^n$.
Thus a cycle has a primitive precisely when its cohomology class is zero.
\end{definition}

The equation $d^2=0$ places every boundary among the cycles, so the quotient
is defined. The two-term complex above has $H^0=R/(t^r)$.
Its degree $-1$ cohomology is zero because $at^r=0$ forces $a=0$ in $R$.
The element $\tau$ is a primitive of $t^r$ in this complex.
For a differential graded algebra, the graded Leibniz rule makes
cohomology an associative graded algebra.
Indeed, cycles multiply to cycles, while $(dx)b=d(xb)$ for a cycle $b$.
For a cycle $a$, the other order satisfies $a(dx)=(-1)^{|a|}d(ax)$.

The scalar cohomology map is $R\to H^0(A)$, $a\mapsto[a1]$.
It is defined because $d(a1)=a\,d1=0$.

\begin{proposition}\label{prop:rel-state-cohomology}
The cohomology algebra of $E$ is
\[
 H^*(E)=k[t,u]/(t^r),\qquad |t|=0,\quad |u|=-2.
\]
Thus $H^{-2j}(E)=R/(t^r)\,[u^j]$ for $j\geq0$, and every other
cohomology group is zero. The kernel of $R\to H^0(E)$ is $(t^r)$.
\end{proposition}

\begin{proof}
Every element of $R[u]$ is a cycle. If $\tau b(u)$ is a cycle, then
$s b(u)=0$. The polynomial ring $k[t,u]$ has no zero divisors, so $b=0$.
The boundaries in the even part are exactly $sR[u]$.
Multiplication of the surviving even words gives the stated algebra.
In degree zero this quotient is $R/(s)$, with the indicated scalar kernel.
\end{proof}

For example, the part in degrees $-3$ through zero is
\[
 R\tau^3\xrightarrow{\ s\ }Ru
 \xrightarrow{\ 0\ }R\tau\xrightarrow{\ s\ }R.
\]
Each odd word is a primitive of $s$ times the next even word.
The class $[u]$ is nonzero for every $r\geq1$.

\section{A homotopy and its multiplication defect}
\label{sec:rel-state-homotopy}

A map of complexes, also called a chain map, is a degree-zero linear map
$F:V\to W$ satisfying $d_WF=Fd_V$.
It sends a cycle $v$ to a cycle because $d_WF(v)=F(d_Vv)=0$.
It sends a boundary $d_Vx$ to the boundary $d_WF(x)$.
Consequently
\[
 H^n(F):H^n(V)\longrightarrow H^n(W),\qquad [v]\longmapsto[F(v)]
\]
is well-defined: changing $v$ by a boundary changes $F(v)$ by a boundary.
The formula preserves identities and composition, since evaluation on
a representative gives $H^n(1)[v]=[v]$ and
$H^n(GF)[v]=[G(F(v))]=H^n(G)H^n(F)[v]$.

Two chain maps $F,G$ are homotopic if a degree-minus-one linear map
$h:V\to W$ satisfies
\[
 F-G=d_Wh+hd_V.
\]
They then induce the same cohomology map: for a cycle $v$, the difference
$F(v)-G(v)$ is the boundary $d_W(hv)$.
A map homotopic to zero is null-homotopic, agreeing with the scalar
definition above.

Left multiplication by $\tau$ defines a degree-minus-one $R$-linear map
$L_\tau:E\to E$. For homogeneous $a$, the Leibniz rule gives
\begin{equation}\label{eq:rel-left-homotopy}
 d(\tau a)+\tau\,da=s a.
\end{equation}
Thus $L_\tau$ is a null-homotopy for scalar multiplication by $s$.
This proves directly that $s$ annihilates every cohomology class of $E$.
The scalar $t^{r-1}$ does not annihilate the unit class, by
Proposition~\ref{prop:rel-state-cohomology}.

The map $L_\tau$ respects the right action, since $\tau(mb)=(\tau m)b$.
Its left compatibility fails at $a=\tau$, $m=b=1$:
\[
 L_\tau(\tau)=u,\qquad
 (-1)^{|\tau|}\tau L_\tau(1)=-u.
\]
The difference is $2u$. Its cohomology class is nonzero because $2$ is
invertible in $k$ and $[u]\ne0$.
Thus the null-homotopy $L_\tau$ fails the left multiplication identity
by the nonzero cohomology class $2[u]$.

There is also no alternative strict two-sided null-homotopy for $s$ on
$E$. Any such map $h$ would satisfy $d(h(1))=s$.
Since $E^{-1}=R\tau$, write $h(1)=a(t)\tau$.
Then $a(t)s=s$ forces $a(t)=1$.
Right compatibility gives $h(\tau)=u$, while left compatibility gives
$h(\tau)=-u$. These equations contradict $2u\ne0$ in $E$.

The earlier equations $dh=\lambda$ and $db=\tau h+h\tau$ use elements
$h\in E^{-1}$ and $b\in E^{-3}$.
The second gives a primitive of the left defect of $L_h$.
It does not make $L_h$ strictly compatible with both actions on $E$.
To incorporate the correction into a homotopy, we will construct a
complex mapping to $E$ with both multiplication actions.
An additional generator will have differential equal to the difference
of the two actions of $\tau$.
On that complex, a scalar null-homotopy respecting both actions will
exist precisely when the two displayed equations have a solution.

\section{Maps from the free algebra}

A morphism of unital differential graded $R$-algebras is an $R$-linear
map that preserves degree, multiplication, and the unit, and commutes
with the differentials.

To compare descriptions of the same algebra, we need inverse maps.
An injective map sends distinct inputs to distinct outputs.
A surjective map attains every output.
A bijection has both properties, so each output has a unique input.
Assigning that input to the output defines its inverse function.

A bijective module map is a module isomorphism because its inverse
preserves sums and scalars: apply the original map to either proposed
identity and use injectivity.
A differential graded algebra isomorphism is a morphism with an inverse
morphism.
For a bijective morphism, decompose an inverse image into homogeneous
components. Injectivity forces every component of the wrong degree to vanish.
Its inverse also preserves products, the unit, and the differential:
apply the original morphism to each required identity and use injectivity.

\begin{proposition}[The universal property of adjoining a primitive]
\label{prop:rel-witness-mapping}
For every unital differential graded $R$-algebra $A$, evaluation at
$\tau$ gives a bijection from its set of unital differential graded
$R$-algebra morphisms $E\to A$ to
\[
 \{h\in A^{-1}:dh=t^r1\}.
\]
Thus specifying a morphism from $E$ is equivalent to choosing a primitive
of $t^r1$, with no additional equation on that primitive.
\end{proposition}

\begin{proof}
Such a morphism sends $\tau$ to an element $h$ of degree $-1$.
Commutation with the differentials gives $dh=t^r1$.
Multiplicativity and $R$-linearity force its value on every element:
\[
 \sum_{i\geq0}a_i\tau^i\longmapsto\sum_{i\geq0}a_i h^i.
\]
All sums are finite, so this formula is defined in $A$.

Conversely, for any such $h$ the displayed formula is $R$-linear,
preserves degree and the unit, and preserves multiplication.
The Leibniz rule gives
\[
 d(h^n)=t^r\sum_{i=0}^{n-1}(-1)^i h^{n-1}.
\]
This agrees with the image of
\eqref{eq:rel-word-differential}, proving that the map commutes with
the differentials. The two constructions are inverse.
\end{proof}

The correspondence respects maps of target algebras.
If $F:A\to B$ is a morphism and $dh=t^r1_A$, then
$dF(h)=F(dh)=t^r1_B$. Composing the evaluation map at $h$ with $F$
therefore gives the evaluation map at $F(h)$.
This compatibility is the meaning of naturality here. Both the algebra map and the induced map on primitives go from
$A$ to $B$.

It also characterizes the universal pair $(E,\tau)$ uniquely.
If another pair $(E',\tau')$ has the same property, evaluation provides
unique morphisms $v:E\to E'$ and $w:E'\to E$ carrying the specified
primitive of each pair to that of the other. Both $wv$ and the identity carry $\tau$ to $\tau$, so
uniqueness gives $wv=1_E$. Likewise $vw=1_{E'}$.
Thus the pair is unique up to the unique isomorphism preserving its
specified primitive.

Imposing $\tau^2=0$ would change this mapping property.
If an ideal $I$ is a sum of homogeneous parts and $dI\subset I$, its
quotient inherits a grading and differential $d[a]=[da]$.
The containment makes the latter independent of the representative,
and the differential identities pass to classes.
For $I=(\tau^2)$, the Leibniz rule and $d(\tau^2)=0$ give $dI\subset I$.
Thus $E/(\tau^2)$ is a differential graded algebra.

Its morphisms to $A$ correspond exactly to primitives
$h$ with both $dh=t^r1$ and $h^2=0$.
Indeed, the word evaluation map factors through that quotient precisely
when it kills $\tau^2$ and hence every multiple of $\tau^2$.

An associative graded algebra need not be graded commutative.
The latter condition would require $ab=(-1)^{|a||b|}ba$ for homogeneous
elements. Applied to $a=b=\tau$, it would give $2\tau^2=0$.
Our algebra imposes no such relation. Although its single-generator
multiplication commutes in the ordinary sense, it is not graded
commutative.


\section{Two actions as one algebra action}
\label{sec:rel-bimodules}

A differential graded left $E$-module is a complex $M$ with a unital
associative left action that adds degrees and satisfies
$d(am)=(da)m+(-1)^{|a|}a\,dm$.
A right module has an action satisfying
$d(ma)=(dm)a+(-1)^{|m|}m\,da$.
An $E$-bimodule has both actions, with $(am)b=a(mb)$.
The left and right actions of each scalar in $R$ agree.
Together the differential rules give
\begin{equation}\label{eq:rel-bimodule-leibniz}
 d(amb)=(da)mb+(-1)^{|a|}a(dm)b
       +(-1)^{|a|+|m|}am(db).
\end{equation}
The algebra $E$ with its multiplication actions is the diagonal $E$-bimodule.

The tensor product $M\otimes_R N$ of $R$-modules consists of finite sums
of symbols $m\otimes n$, subject to additivity and scalar balancing:
\[
 (am)\otimes n=m\otimes(an)=a(m\otimes n),\qquad a\in R.
\]
For graded modules the degree of $m\otimes n$ is $|m|+|n|$.
For complexes its differential is
\begin{equation}\label{eq:rel-tensor-differential}
 d(m\otimes n)=dm\otimes n+(-1)^{|m|}m\otimes dn.
\end{equation}
Applying $d$ twice cancels the two terms containing $dm\otimes dn$.

For homogeneous linear maps $F,G$, define
\[
 (F\otimes G)(x\otimes y)=(-1)^{|G||x|}F(x)\otimes G(y).
\]
Additivity and scalar balancing make this a map of degree $|F|+|G|$.
For degree-zero chain maps, the two terms of the tensor differential
show directly that $F\otimes G$ is a chain map.

Repeated tensor products have the reassociation isomorphism
\[
 (M\otimes_R N)\otimes_R Q\longrightarrow
 M\otimes_R(N\otimes_R Q),\qquad
 (m\otimes n)\otimes q\longmapsto m\otimes(n\otimes q).
\]
The formula in reverse defines its inverse. Both respect the defining
additivity and balancing relations. The differential in either grouping
is the sum with differentiated entries $m,n,q$ and respective signs
$1,(-1)^{|m|},(-1)^{|m|+|n|}$.
We therefore omit parentheses when writing three tensor factors.

For graded algebras the product is
\[
 (a\otimes b)(c\otimes e)=(-1)^{|b||c|}ac\otimes be.
\]
This sign records the interchange of $b$ and $c$ in the two tensor factors.

When $M$ is a right $E$-module and $N$ is a left $E$-module, define
$M\otimes_E N$ by the additional relations
\[
 (ma)\otimes n=m\otimes(an).
\]
There is no sign in this balancing relation, which does not interchange
the two inputs. Formula~\eqref{eq:rel-tensor-differential} respects the
relation. In fact, its two sides have the same three terms containing
$dm$, $da$, and $dn$, by the module Leibniz rules.
If both inputs are bimodules, the remaining outer actions make this tensor
product a bimodule.

Reassociation is defined by the same formula for balanced tensor products
over $E$, since moving an element of $E$ between adjacent factors gives
the same relation in either grouping.
For bimodule maps, the tensor-map formula also respects this balancing:
the values on $(xa)\otimes y$ and $x\otimes ay$ both have sign
$(-1)^{|G|(|x|+|a|)}$ after graded left linearity of $G$ is applied.
Here graded left linearity means
$G(ay)=(-1)^{|G||a|}aG(y)$, whereas right linearity means
$F(xa)=F(x)a$. These are precisely the compatibility conventions we will
derive from the enveloping action.

The action map $E\otimes_E M\to M$ sends $a\otimes m$ to $am$.
Its inverse is $m\mapsto1\otimes m$, because balancing gives
$1\otimes am=a\otimes m$.
Similarly, $M\otimes_E E\to M$, $m\otimes a\mapsto ma$, has inverse
$m\mapsto m\otimes1$.
The module differential rules show that these are chain maps.

Define the graded opposite algebra $E^{\mathrm{op}}$ to have the same
complex as $E$, with multiplication
\[
 a^{\mathrm{op}}b^{\mathrm{op}}
   =(-1)^{|a||b|}(ba)^{\mathrm{op}}.
\]
The sign ensures the graded Leibniz rule: both
$d(a^{\mathrm{op}}b^{\mathrm{op}})$ and
$(da)^{\mathrm{op}}b^{\mathrm{op}}
 +(-1)^{|a|}a^{\mathrm{op}}(db)^{\mathrm{op}}$ equal
\[
 (-1)^{|a||b|}\bigl((db)a\bigr)^{\mathrm{op}}
 +(-1)^{|a||b|+|b|}\bigl(b(da)\bigr)^{\mathrm{op}}.
\]
The enveloping algebra is
\[
 E^e=E\otimes_R E^{\mathrm{op}}.
\]
An $E$-bimodule becomes a left $E^e$-module by
\begin{equation}\label{eq:rel-envelope-action}
 (a\otimes b^{\mathrm{op}})m=(-1)^{|b||m|}amb.
\end{equation}

To check associativity, let $a,b,c,e,m$ be homogeneous.
Acting first by $c\otimes e^{\mathrm{op}}$ and then by
$a\otimes b^{\mathrm{op}}$ gives the word $acmeb$ with sign exponent
\[
 |e||m|+|b|(|c|+|m|+|e|).
\]
Multiplying inside $E^e$ first gives the same word with exponent
\[
 |b||c|+|b||e|+(|b|+|e|)|m|.
\]
These exponents agree.

For $x=a\otimes b^{\mathrm{op}}$, both
$d(xm)$ and $(dx)m+(-1)^{|x|}x(dm)$ expand to
\[
 (-1)^{|b||m|}\bigl((da)mb+(-1)^{|a|}a(dm)b
                +(-1)^{|a|+|m|}am(db)\bigr).
\]
This proves differential compatibility.
The actions of $a\otimes1$ and $1\otimes b^{\mathrm{op}}$ recover the
original two actions, so no action data have been discarded.

A homogeneous linear map of degree $n$ between graded left modules is
graded module-linear when it additionally satisfies
$F(xm)=(-1)^{n|x|}xF(m)$ for homogeneous $x$.
Using \eqref{eq:rel-envelope-action}, this is equivalent for bimodules to
\begin{equation}\label{eq:rel-graded-bimodule-map}
 F(amb)=(-1)^{n|a|}aF(m)b.
\end{equation}
For the right action, the signs $(-1)^{|b||m|}$ on the input and
$(-1)^{|b|(|m|+n)}$ on the output cancel the graded map sign.
For the left action only $(-1)^{n|a|}$ remains.
This proves the convention used in
\eqref{eq:rel-homotopy-linearity}.

Write $\operatorname{Hom}_{E-E}(M,N)^n$ for the $R$-module of additive
$R$-linear maps of degree $n$ satisfying
\eqref{eq:rel-graded-bimodule-map}, with pointwise addition and scalar
multiplication. A map may have nonzero values in
infinitely many input degrees. Its output on each individual element is
an element of the algebraic module $N$.
The Hom differential is
\begin{equation}\label{eq:rel-hom-differential}
 \partial F=d_NF-(-1)^nFd_M.
\end{equation}

The bimodule Leibniz rule cancels the terms containing $da$ and $db$ in
the Hom differential, giving
\[
 (\partial F)(amb)=(-1)^{(n+1)|a|}a(\partial F)(m)b.
\]
Thus $\partial F$ has the required linearity in degree $n+1$.
Expanding twice gives $\partial^2F=0$.
The degree-zero cycles are bimodule chain maps, and their boundaries are
the bimodule maps null-homotopic through compatible maps.

An ordinary degree-zero bimodule endomorphism $F:E\to E$ is determined
by $c=F(1)$. Its two action identities give $F(a)=ac=ca$.
If it is a chain map, then $dc=0$ as well.
The ordinary center therefore imposes commutation as an equality.
The derived center will instead use a free bimodule resolution of the
diagonal, with additional generators whose differentials impose the
action relations.

\section{An explicit two-cell resolution}
\label{sec:rel-resolution}

A quasi-isomorphism is a chain map inducing an isomorphism on every
cohomology group. We now construct one from a free graded bimodule to $E$.
Put
\[
 P=Ee_0E\oplus Ee_1E,
 \qquad Ee_cE=E\otimes_R Re_c\otimes_R E,
 \qquad |e_0|=0,\quad |e_1|=-2.
\]
The notation $ae_cb$ denotes $a\otimes e_c\otimes b$.
There are two free generators, or cells, in these specified degrees.

A graded bimodule map $F:P\to N$ of degree $n$ is determined by its
values on these generators, through
\[
 F(ae_cb)=(-1)^{n|a|}aF(e_c)b.
\]
Conversely, any values of degrees $n$ and $n-2$ extend by this formula
and additivity to such a map. Balancing over $R$ makes the extension
well-defined.

Define
\begin{equation}\label{eq:rel-resolution-differential}
 d_Pe_0=0,\qquad d_Pe_1=\tau e_0-e_0\tau,
\end{equation}
and extend using \eqref{eq:rel-bimodule-leibniz}.
The output on $e_1$ has degree $-1$, as required.
Its next differential is $s e_0-e_0s=0$, since the tensor products are
over $R$. The graded Leibniz rule then gives $d_P^2=0$ everywhere.

Define the augmentation and its $R$-linear section by
\[
 \begin{split}
 q:P\longrightarrow E,&\qquad q(ae_0b)=ab,\quad q(ae_1b)=0,\\
 \iota:E\longrightarrow P,&\qquad \iota(a)=e_0a.
 \end{split}
\]
Both are chain maps, and $q\iota=\operatorname{id}_E$.
The map $q$ is a bimodule map. The section $\iota$ is right $E$-linear
but need not be left $E$-linear.

\begin{proposition}\label{prop:rel-contraction}
The augmentation $q$ is a quasi-isomorphism. More precisely, the degree
minus one $R$-linear map
\begin{equation}\label{eq:rel-contraction}
 \begin{split}
 h(\tau^i e_0\tau^j)
 &=\sum_{\ell=0}^{i-1}(-1)^\ell
       \tau^\ell e_1\tau^{i+j-1-\ell},\\
 h(\tau^i e_1\tau^j)&=0
 \end{split}
\end{equation}
satisfies $d_Ph+hd_P=\operatorname{id}_P-\iota q$ for all $i,j\geq0$.
\end{proposition}

\begin{proof}
An empty sum is zero. Each term in the first formula has degree
$-(i+j+1)$, so $h$ lowers degree by one.
Write $X_{ij}=\tau^i e_0\tau^j$ and $Y_{ij}=\tau^i e_1\tau^j$.
Set $\epsilon_i=0$ for even $i$ and $\epsilon_i=1$ for odd $i$.

We split $d_P=s\delta+\rho$, where
\[
 \begin{split}
 \delta X_{ij}&=\epsilon_iX_{i-1,j}
                       +(-1)^i\epsilon_jX_{i,j-1},\\
 \delta Y_{ij}&=\epsilon_iY_{i-1,j}
                       +(-1)^i\epsilon_jY_{i,j-1},\\
 \rho X_{ij}&=0,\\
 \rho Y_{ij}&=(-1)^i(X_{i+1,j}-X_{i,j+1}).
 \end{split}
\]
A term with a negative index and zero coefficient is omitted.
The factor $(-1)^i$ in $\rho$ comes from passing the differential across
the left word $\tau^i$.

On $X_{ij}$ the relation terms telescope:
\[
 \rho hX_{ij}
 =\sum_{\ell=0}^{i-1}
     \bigl(X_{\ell+1,i+j-1-\ell}-X_{\ell,i+j-\ell}\bigr)
 =X_{ij}-X_{0,i+j}.
\]
On $Y_{ij}$ they give
\[
 h\rho Y_{ij}
 =(-1)^i\bigl(hX_{i+1,j}-hX_{i,j+1}\bigr)=Y_{ij}.
\]
The two sums in parentheses have the same first $i$ terms.
Their remaining term is $(-1)^iY_{ij}$.
Thus $\rho h+h\rho=\operatorname{id}_P-\iota q$.

It remains to prove $\delta h+h\delta=0$.
Both summands vanish on every $Y_{ij}$.
They also vanish on $X_{0j}$, since $hX_{0j}=0$ and $h\delta X_{0j}=0$.

For $i\geq1$, put $N=i+j-1$. For $0\leq p\leq i-2$, the coefficient
of $Y_{p,N-1-p}$ in $(\delta h+h\delta)X_{ij}$ is
\[
 \epsilon_{N-p}
 +(-1)^p\bigl(-\epsilon_{p+1}+\epsilon_i+(-1)^i\epsilon_j\bigr).
\]
Use the identities
\[
 \epsilon_{a+b}=\epsilon_a+(-1)^a\epsilon_b,
 \qquad \epsilon_{p+1}+(-1)^{p+1}\epsilon_{N-p}
       =\epsilon_{i+j}.
\]
The coefficient is zero because
$\epsilon_i+(-1)^i\epsilon_j=\epsilon_{i+j}$.

If $j\geq1$, the last possible coefficient, at $p=i-1$, is
\[
 \epsilon_j+(-1)^i\epsilon_j(-1)^{i-1}=0.
\]
For $j=0$ this last basis term does not occur.
These are all terms, proving the required cancellation.

The two identities for $\rho$ and $\delta$ prove the contraction formula.
It implies that $q$ and $\iota$ induce inverse cohomology maps.
Consequently $q$ is a quasi-isomorphism.
\end{proof}

The contraction has a direct interpretation in its first cases:
\[
 h(\tau e_0)=e_1,
 \qquad h(\tau^2e_0)=e_1\tau-\tau e_1.
\]
For the first word, $d_Ph(\tau e_0)=\tau e_0-e_0\tau$.
For the second, differentiating the two terms gives opposite scalar
contributions $s e_1$ and $-s e_1$.
Their remaining terms equal $\tau^2e_0-e_0\tau^2$.
The general formula extends these two cancellations to every word.

\section{Why this resolution computes derived maps}
\label{sec:rel-derived-maps}

To compute maps from a resolution, we need its Hom complex to preserve
quasi-isomorphisms in the target.
A bimodule $Q$ is homotopically projective if
$\operatorname{Hom}_{E-E}(Q,N)$ is acyclic for every acyclic bimodule $N$.
A bimodule quasi-isomorphism $q:Q\to E$ from such a $Q$ is a
homotopically projective resolution of the diagonal.

\begin{lemma}\label{lem:rel-hom-acyclic}
If the bimodule $N$ is acyclic, then
$\operatorname{Hom}_{E-E}(P,N)$ is acyclic.
\end{lemma}

\begin{proof}
A degree-$n$ map $F:P\to N$ is determined by
\[
 F(e_0)=a\in N^n,\qquad F(e_1)=b\in N^{n-2}.
\]
Its differential, evaluated on the two generators, is
\begin{equation}\label{eq:rel-pair-differential}
 \partial(a,b)=\bigl(da,\,db-\tau a+(-1)^n a\tau\bigr).
\end{equation}
For the second component, graded left linearity gives
$F(\tau e_0)=(-1)^n\tau a$, while right linearity gives
$F(e_0\tau)=a\tau$. Substitution into
\eqref{eq:rel-hom-differential} proves the formula.

Suppose $(a,b)$ is a cycle. Since $da=0$, acyclicity supplies
$\alpha\in N^{n-1}$ with $d\alpha=a$.
Subtracting $\partial(\alpha,0)$ leaves a cycle $(0,\beta)$.
Its second component satisfies $d\beta=0$ by
\eqref{eq:rel-pair-differential}.

Choose $\gamma\in N^{n-3}$ with $d\gamma=\beta$.
Then $(0,\beta)=\partial(0,\gamma)$, so
$(a,b)=\partial(\alpha,\gamma)$.
Every cycle is a boundary.
\end{proof}

The cone construction relates acyclic targets to quasi-isomorphisms.
For a chain map $f:V\to W$, its cone is the
complex with
\[
 \operatorname{Cone}(f)^n=W^n\oplus V^{n+1},\qquad
 d(b,a)=(db+f(a),-da).
\]
The second component has had its degree decreased by one.
The chain map equation implies that the displayed differential squares to
zero.

For a degree-zero bimodule chain map $f:V\to W$, postcomposition gives
\[
 \operatorname{Hom}_{E-E}(Q,f):
 \operatorname{Hom}_{E-E}(Q,V)\longrightarrow
 \operatorname{Hom}_{E-E}(Q,W),\qquad F\longmapsto f\circ F.
\]
For a map $F$ of degree $n$, the chain map equation for $f$ gives
\[
 \partial(fF)=d_WfF-(-1)^n fF d_Q
 =f\bigl(d_VF-(-1)^nF d_Q\bigr)=f\,\partial F.
\]
Thus postcomposition is a chain map between the two Hom complexes.

\begin{lemma}\label{lem:rel-cone}
The cone of $f$ is acyclic if and only if $f$ is a quasi-isomorphism.
If $Q$ has the acyclic-target property in the definition above, then
$\operatorname{Hom}_{E-E}(Q,f)$ is a quasi-isomorphism for every bimodule
quasi-isomorphism $f$.
\end{lemma}

\begin{proof}
First suppose $f$ induces cohomology isomorphisms.
A cone cycle $(b,a)$ satisfies $da=0$ and $f(a)=-db$.
Injectivity gives $a=dx$ for some $x$.
Then $b+f(x)$ is a cycle, so surjectivity supplies a cycle $c$ and an
element $y$ with $b+f(x)=f(c)+dy$.
It follows that $(b,a)=d(y,c-x)$.

Conversely, suppose the cone is acyclic.
For a cycle $b\in W^n$, write $(b,0)=d(y,x)$.
Then $dx=0$ and $b=dy+f(x)$, proving surjectivity on cohomology.
If $a\in V^n$ is a cycle with $f(a)=db$, then $(-b,a)$ is a cone cycle.
A cone primitive has second component $x$ with $a=-dx$.
This proves injectivity.

For bimodules, give the cone the actions
\[
 \alpha(b,a)\beta
 =\bigl(\alpha b\beta,(-1)^{|\alpha|}\alpha a\beta\bigr).
\]
The left sign compensates for the degree decrease in the second component.
Expanding the differential on these two components verifies
\eqref{eq:rel-bimodule-leibniz}.

A degree-$n$ map into this cone is a pair $(F,G)$ of bimodule maps
of degrees $n,n+1$ into $W,V$.
Its differential is $(\partial F+fG,-\partial G)$.
Thus there is an equality of complexes under component evaluation:
\[
 \operatorname{Hom}_{E-E}(Q,\operatorname{Cone}(f))
 =\operatorname{Cone}\bigl(\operatorname{Hom}_{E-E}(Q,f)\bigr).
\]

The first part makes $\operatorname{Cone}(f)$ acyclic.
The acyclic-target property makes the left side acyclic.
The first part applied again proves the second assertion.
\end{proof}

Proposition~\ref{prop:rel-contraction} and
Lemma~\ref{lem:rel-hom-acyclic} show that $q:P\to E$ is a
homotopically projective resolution.
For any bimodule $N$, the groups
$H^n\operatorname{Hom}_{E-E}(P,N)$ are the derived bimodule maps from the
diagonal to $N$. The word derived specifies that we compute from such a
resolution. Lemma~\ref{lem:rel-cone} proves that these groups respect
quasi-isomorphisms in the target.

The resolution choice also does not affect the resulting maps or their
composition algebra. To see this, let $q':Q\to E$ be another
homotopically projective resolution. The quasi-isomorphism
\[
 \operatorname{Hom}_{E-E}(P,Q)\xrightarrow{\ q'\circ(-)\ }
 \operatorname{Hom}_{E-E}(P,E)
\]
provides a degree-zero chain map $f:P\to Q$ with $q'f$ homotopic to $q$.
Interchanging $P,Q$ provides $g:Q\to P$ with $qg$ homotopic to $q'$.

The maps $q(gf)$ and $q$ are homotopic.
Injectivity on degree-zero cohomology of $\operatorname{Hom}(P,q)$ gives
$gf\simeq\operatorname{id}_P$.
The same argument gives $fg\simeq\operatorname{id}_Q$.

Here and below composition is written $FG=F\circ G$, so $G$ acts first.
Write $\operatorname{End}_{E-E}(M)=\operatorname{Hom}_{E-E}(M,M)$
for maps from a bimodule to itself.
For homogeneous maps the Hom differential satisfies
\[
 \partial(FG)=(\partial F)G+(-1)^{|F|}F(\partial G).
\]
This follows by expanding both sides using
\eqref{eq:rel-hom-differential}.
It proves that composition descends to cohomology.

The maps
\[
 [F]\longmapsto[fFg],\qquad [G]\longmapsto[gGf]
\]
then give inverse graded algebra isomorphisms between
$H^*\operatorname{End}_{E-E}(P)$ and
$H^*\operatorname{End}_{E-E}(Q)$.
Indeed, in cohomology $[gf]=[\operatorname{id}_P]$ and
$[fg]=[\operatorname{id}_Q]$.
These two identities prove both preservation of products and the inverse
property. Different lifts $f,g$ are homotopic by the same injectivity
argument, so the induced isomorphism is independent of those lifts.

We define the relative Hochschild cohomology, or cohomology of the
relative derived center, by
\[
 \operatorname{HH}_R^*(E)
   =H^*\operatorname{End}_{E-E}(P).
\]
The endomorphism complex has composition as product and identity map as
unit. The subscript $R$ records the common scalar action in the two tensor
factors of $E^e$. The preceding argument makes this definition independent
of the homotopically projective resolution.

The scalar relation now gives a precise comparison problem.
Multiplication by $t^r$ is null-homotopic as an $R$-linear map on $E$.
We ask which powers of $t$ admit a bimodule null-homotopy on $P$.
Equivalently, we must compute the kernel of
$R\to\operatorname{HH}_R^0(E)$.
The calculation below gives $(t^{2r})$, whereas the kernel of
$R\to H^0(E)$ is $(t^r)$.
For $r=2$, the first bimodule null-homotopy therefore occurs at $t^4$.

Tensoring over $R$ with $E$ preserves quasi-isomorphisms.
As an $R$-complex, $E$ is a direct sum, for $j\geq0$, of
\[
 R\tau u^j\xrightarrow{\ s\ }Ru^j
\]
in degrees $-2j-1,-2j$.

We prove this first for acyclic targets.
A subcomplex $U\subset V$ consists of submodules $U^n\subset V^n$
with $d(U^n)\subset U^{n+1}$.
The quotient modules $V^n/U^n$ form a complex with differential
$[v]\mapsto[dv]$, well-defined by this containment.

Tensor one displayed summand with an acyclic $R$-complex $N$.
Its subcomplex $Ru^j\otimes_R N$ is a copy of $N$ with degrees translated
by $-2j$. The quotient is a copy of $N$ with degrees translated by
$-2j-1$ and differential $-d_N$.
Both are acyclic.

Here a short exact sequence of complexes is degreewise an injection,
followed by a surjection whose kernel is the image of that injection.
In such a sequence with acyclic subcomplex and quotient, the middle
complex is acyclic.
For a proof, send a middle cycle to the quotient and choose a primitive
there. Subtract the differential of a lift of that primitive.
The remaining cycle belongs to the subcomplex and has a primitive there.

This proves the assertion for each two-term tensor summand.
Their direct sum is acyclic because every element has finite support.
Thus tensoring with $E$ preserves acyclicity.

To apply the cone criterion, for an $R$-complex $A$ use the isomorphism
\[
 A\otimes_R\operatorname{Cone}(f)\longrightarrow
 \operatorname{Cone}(1_A\otimes f),\qquad
 a\otimes(w,v)\longmapsto
 \bigl(a\otimes w,(-1)^{|a|}a\otimes v\bigr).
\]
It is its own inverse on components. Under either route through the
differentials, the first component is
\[
 da\otimes w+(-1)^{|a|}a\otimes dw+(-1)^{|a|}a\otimes f(v),
\]
and the second is $(-1)^{|a|+1}da\otimes v-a\otimes dv$.
It is therefore an isomorphism of complexes.
Taking $A=E$ and applying the cone criterion proves preservation of
quasi-isomorphisms by $E\otimes_R(-)$.

The interchange map $a\otimes v\mapsto(-1)^{|a||v|}v\otimes a$
is an isomorphism of $R$-complexes: applying it twice gives the identity,
and the tensor differentials agree, with signs
$(-1)^{(|a|+1)|v|}$ on $v\otimes da$ and
$(-1)^{|a||v|}$ on $dv\otimes a$.
Hence tensoring on the other side preserves quasi-isomorphisms as well.
The same statement holds for $E^{\mathrm{op}}$, whose underlying complex
is unchanged. Hence any quasi-isomorphism $V\to E$ of $R$-complexes
induces a quasi-isomorphism
\[
 V\otimes_R E^{\mathrm{op}}\longrightarrow
 E\otimes_R E^{\mathrm{op}}.
\]
The corresponding assertion holds with the first factor fixed.
This is the tensor invariance required to use the ordinary enveloping
algebra in this relative calculation.

\section{The central complex and its product}
\label{sec:rel-central-complex}

The complex $C=\operatorname{Hom}_{E-E}(P,E)$ is smaller than the
endomorphism complex of $P$. We first compute its differential and then
give it the composition product through an explicit comparison.
A degree-$n$ element is a pair $(a,b)$ with
\[
 a\in E^n,\qquad b\in E^{n-2}.
\]
In particular, the second component has total degree $|b|+2$ as an
element of $C$. Formula~\eqref{eq:rel-pair-differential} gives
\begin{equation}\label{eq:rel-central-pair}
 D(a,b)=\bigl(da,\,db-\tau a+(-1)^n a\tau\bigr).
\end{equation}
Here $D$ denotes the Hom differential, while $d$ remains the differential
on $E$.

The commutation term cancels in $D^2$ as follows.
For $a\in E^n$, write $K_n(a)=-\tau a+(-1)^n a\tau$.
Expansion with $d\tau=s$ gives
\[
 dK_n(a)+K_{n+1}(da)=-sa+as=0.
\]
The first component of $D^2$ is $d^2a=0$.
The second is $d^2b+dK_n(a)+K_{n+1}(da)=0$.
The final equality uses the common scalar action of $s$.

Define a degree-zero bimodule map
\[
 \Delta:P\longrightarrow P\otimes_E P,
 \qquad
 \Delta(e_0)=e_0\otimes e_0,
 \qquad
 \Delta(e_1)=e_1\otimes e_0+e_0\otimes e_1.
\]
The tensor product uses the right action on the first $P$ and the left
action on the second $P$. Its outer actions extend the displayed formulas
uniquely to a bimodule map.

\begin{lemma}\label{lem:rel-diagonal}
The map $\Delta$ is a chain map. It is coassociative and has counit $q$:
\[
 (\Delta\otimes1)\Delta=(1\otimes\Delta)\Delta,
 \qquad (q\otimes1)\Delta=(1\otimes q)\Delta=\operatorname{id}_P,
\]
using the action identifications $E\otimes_E P=P=P\otimes_E E$.
\end{lemma}

\begin{proof}
Both sides of the chain map equation vanish on $e_0$.
On $e_1$, the differential of its image is
\[
 \begin{split}
 d\Delta(e_1)
  ={}&\tau e_0\otimes e_0-e_0\tau\otimes e_0\\
    &+e_0\otimes\tau e_0-e_0\otimes e_0\tau\\
  ={}&\tau e_0\otimes e_0-e_0\otimes e_0\tau
   =\Delta(d_Pe_1).
 \end{split}
\]
The middle terms cancel by balancing.

Both iterated diagonals send $e_0$ to $e_0\otimes e_0\otimes e_0$.
Both send $e_1$ to the sum of three tensors with $e_1$ in one position
and $e_0$ in the other two positions.
This proves coassociativity on the generators and hence on $P$.
Finally, $q(e_0)=1$ and $q(e_1)=0$ give both counit identities on the
generators.
\end{proof}

Let $\mu:E\otimes_E E\to E$ be multiplication.
Define the cup product by
\[
 F\smile G=\mu(F\otimes G)\Delta.
\]
Both $e_0$ and $e_1$ have even degree. Thus evaluation on the generators
introduces no further tensor signs, and the formula is
\begin{equation}\label{eq:rel-cup-product}
 (a,b)\smile(c,e)=(ac,bc+ae).
\end{equation}
This formula is associative by direct multiplication of three pairs.
Its unit is $(1,0)$, which is the augmentation $q$ as an element of $C$.

Introduce $z=(0,1)$, a homogeneous element of degree two.
Formula~\eqref{eq:rel-cup-product} gives $z^2=0$ and
$z(a,0)=(a,0)z=(0,a)$.
Consequently the graded cup algebra is
\[
 C=E[z]/(z^2),\qquad a+bz\longleftrightarrow(a,b),\qquad |z|=2,
\]
where $z$ commutes with every element of $E$.
The notation is a quotient of an associative graded algebra, and still
retains $u=\tau^2$.

Formula~\eqref{eq:rel-central-pair} now reads
\begin{equation}\label{eq:rel-central-model}
 Dt=0,\qquad Dz=0,\qquad Du=0,\qquad D\tau=s-2uz.
\end{equation}
For $\tau$ the second component is
$-\tau\tau-\tau\tau=-2u$.
For $u$ it is $-\tau u+u\tau=0$.
Thus the coefficient $2$ has been determined by graded linearity and the
specified generator $e_1$.

These formulas define a derivation of the cup product.
For an explicit check, the even subalgebra $R[u,z]/(z^2)$ is killed by
$D$, and each element of $C$ is a sum of an even element and $\tau$
times an even element. Set $f=s-2uz$.
The Leibniz rule holds for products with an even factor because $f$ is
central. For two odd factors $\tau a,\tau b$, with $a,b$ even, it gives
$fa\tau b-\tau afb=0$, agreeing with $D(uab)=0$.
This checks all products by bilinearity.

\begin{proposition}\label{prop:rel-composition}
The differential graded algebra $C$ computes the composition product on
the relative derived center. More precisely, there is a unital
multiplicative quasi-isomorphism
\[
 T:C\longrightarrow\operatorname{End}_{E-E}(P)
\]
defined, for $F=(a,b)$ of degree $n$, by
\[
 T_F(e_0)=e_0a,\qquad T_F(e_1)=e_1a+e_0b.
\]
\end{proposition}

\begin{proof}
Both displayed values have the required degrees.
Extend them to $P$ by graded bimodule linearity.
For $G=(c,e)$, right linearity gives
\[
 \begin{split}
 T_FT_G(e_0)&=e_0ac,\\
 T_FT_G(e_1)&=e_1ac+e_0(bc+ae).
 \end{split}
\]
These are $T_{F\smile G}$ by
\eqref{eq:rel-cup-product}. Also $T_{(1,0)}=\operatorname{id}_P$.

The endomorphism differential evaluated on $e_0$ is $e_0da$.
On $e_1$ it is
\[
 \begin{split}
 (\partial T_F)(e_1)
  ={}&(\tau e_0-e_0\tau)a+e_1da+e_0db\\
     &-(-1)^n\bigl((-1)^n\tau e_0a-e_0a\tau\bigr)\\
  ={}&e_1da+e_0\bigl(db-\tau a+(-1)^n a\tau\bigr).
 \end{split}
\]
The sign on $T_F(\tau e_0)$ is $(-1)^n$.
The resulting expression proves $\partial T_F=T_{DF}$ on both cells.

Postcomposition with $q$ defines a chain map
\[
 q_*:\operatorname{End}_{E-E}(P)\longrightarrow
           \operatorname{Hom}_{E-E}(P,E)=C.
\]
Since $q$ is a quasi-isomorphism, Lemmas~\ref{lem:rel-hom-acyclic}
and~\ref{lem:rel-cone} make $q_*$ a quasi-isomorphism.

Evaluation on the two cells gives $q_*T_F=F$.
On cohomology, $H(q_*)H(T)$ is therefore the identity.
Since $H(q_*)$ is an isomorphism, so is $H(T)$.
The product identity already proved makes this an isomorphism of graded
algebras. No multiplicativity of $q_*$ is needed.
\end{proof}

The comparison identifies the scalar $a\in R$ in $C$ with the scalar
endomorphism $a\operatorname{id}_P$.
Consequently the scalar calculation in $C$ determines the vanishing of
the corresponding derived central operation.

\section{The exact scalar exponent and the full cohomology ring}
\label{sec:rel-center-cohomology}

Let
\[
 S=R[u,z]/(z^2),\qquad |u|=-2,\quad |z|=2,
 \qquad f=s-2uz.
\]
All elements of $S$ have even degree, and $f$ has degree zero.
The unique even--odd decomposition of the central complex is
\[
 C=S\oplus\tau S,\qquad D(a+\tau b)=fb.
\]

The differential thus reduces to multiplication by one element of $S$.
Although $z$ is a zero divisor in $S$, the element $f$ is not.
Indeed, write $b=b_0+zb_1$ with $b_0,b_1\in R[u]$.
Then
\begin{equation}\label{eq:rel-injectivity}
 fb=s b_0+z(s b_1-2u b_0).
\end{equation}
If this is zero, its first coefficient gives $b_0=0$, since $s$ is not
a zero divisor in $R[u]$. The second then gives $b_1=0$.

The ring quotient will follow from this injectivity. To record its
individual degrees, we also need coordinates for linear maps.
An $R$-linear map between finite free modules with bases
$v_1,\ldots,v_m$ and $w_1,\ldots,w_n$ is represented by its matrix
$A=(a_{ij})$, where
\[
 F(v_j)=\sum_{i=1}^n a_{ij}w_i.
\]
Thus column $j$ records the image of the $j$th basis element.
If another map has matrix $B=(b_{ki})$, their composite has matrix
$BA$, since the coefficient of its $k$th output basis element on $v_j$
is $\sum_i b_{ki}a_{ij}$.

The identity matrix has entries $1$ on its diagonal and $0$ elsewhere.
A square matrix is invertible when another matrix multiplies with it
in either order to give the identity. The corresponding linear maps
are then inverse.

When the scalar ring is a field, its modules are called vector spaces.
The span of vectors is the set of their finite linear combinations.
They are linearly independent when a finite linear combination is zero
only if all its coefficients are zero.
A basis is an independent spanning family, equivalently a family giving
each vector a unique finite coordinate expression.
The span is a subspace, meaning a submodule over the field.

For a space with a finite basis, the number of basis vectors is its
dimension. This number does not depend on the basis.
Indeed, start with a spanning list of $n$ vectors and an independent list
$v_1,\ldots,v_m$. Express $v_1$ in the spanning list and replace an entry
with nonzero coefficient by $v_1$. Division by that coefficient shows
that the new list still spans.

After inserting $v_1,\ldots,v_j$, the expression for $v_{j+1}$ has a
nonzero coefficient on an entry not yet replaced, since otherwise the
$v_i$ would be dependent. Replacing that entry continues the process.
There can be at most $n$ replacements, so $m\leq n$.
Applying this inequality to two finite bases in both orders proves
equality of their lengths. A finite spanning list always contains a
basis: repeatedly remove an entry expressible in the others until the
remaining list is independent.

\begin{theorem}\label{thm:rel-full-center}
For the differential graded algebra
$E=k[t]\langle\tau\rangle$ with $|\tau|=-1$ and $d\tau=t^r$,
the relative derived center has cohomology
\begin{equation}\label{eq:rel-center-ring}
 \operatorname{HH}_R^*(E)
  =\frac{k[t,u,z]}{(z^2,t^r-2uz)},
 \qquad |t|=0,\quad |u|=-2,\quad |z|=2.
\end{equation}
Its homogeneous groups are
\[
 \begin{split}
 \operatorname{HH}_R^{-2j}(E)&=R/(t^{2r})\,[u^j],\qquad j\geq0,\\
 \operatorname{HH}_R^2(E)&=R/(t^r)\,[z].
 \end{split}
\]
All other groups vanish. The kernel of $R\to\operatorname{HH}_R^0(E)$
is exactly $(t^{2r})$.
\end{theorem}

\begin{proof}
Equation~\eqref{eq:rel-injectivity} proves that the odd part has no
cycles. The even part consists of cycles, with boundaries $fS$.
Thus $H^*(C)=S/(f)$ as an algebra.
Proposition~\ref{prop:rel-composition} identifies it with the stated
relative derived center.

To determine the homogeneous groups and the scalar kernel, note that
$u^a z^b$ has degree $-2a+2b$, with $a\geq0$ and $b=0,1$.
For $j\geq0$, an $R$-basis of $C^{-2j}$ is
$u^j,u^{j+1}z$. An $R$-basis of $C^{-2j-1}$ is
$\tau u^j,\tau u^{j+1}z$.
The differential between them has matrix
\begin{equation}\label{eq:rel-scalar-matrix}
 \begin{pmatrix}s&0\\-2&s\end{pmatrix}.
\end{equation}
For example, the second column follows from
$D(\tau u^{j+1}z)=(s-2uz)u^{j+1}z=su^{j+1}z$.

Let $v=[u^j]$ and $w=[u^{j+1}z]$.
The two columns impose $sv-2w=0$ and $sw=0$.
Since $2$ is invertible, every class is a multiple of $v$, and these
relations imply $s^2v=0$.

There are no further relations: if $(g,0)$ lies in the matrix image, then
for some $A,B\in R$,
\[
 g=sA,\qquad 0=-2A+sB.
\]
Hence $A=sB/2$ and $g=s^2B/2$.
Conversely, every multiple of $s^2$ occurs by choosing $A=sB/2$.
The quotient is therefore exactly $R/(s^2)$ on $v$.
For $j=0$ the generator is the unit, proving the scalar kernel assertion.

In degree two, the only basis element is $z$.
Its boundaries are the multiples of $D(\tau z)=sz$, giving $R/(s)$.
There are no even elements in degrees greater than two.
The odd groups already vanished by injectivity.
\end{proof}

For a unital graded algebra $B$ with an $R$-action, the kernel of
$R\to B^0$ is also the scalar annihilator of all of $B$.
If a scalar kills every element it kills the unit.
If it kills the unit, multiplication shows that it kills every element.

Define the scalar annihilator ideals of $H^*(E)$ and
$\operatorname{HH}_R^*(E)$ by
\[
 I_{\mathrm{st}}=\ker\bigl(R\longrightarrow H^0(E)\bigr),\qquad
 I_{\mathrm{cen}}=\ker\bigl(R\longrightarrow
                                  \operatorname{HH}_R^0(E)\bigr).
\]
For ideals $I,J\subset R$, their product $IJ$ consists of finite sums
of products $ab$ with $a\in I$ and $b\in J$. Write $I^2=II$.
The preceding computations give
\[
 I_{\mathrm{st}}=(t^r),\qquad
 I_{\mathrm{cen}}=(t^{2r})=I_{\mathrm{st}}^2
                  \subsetneq I_{\mathrm{st}}.
\]
The containment is strict since $r\geq1$.
It concerns the specified parameter ring $R$ and the relative center
computed over that ring.

The central boundary $s^2$ has an explicit primitive:
\begin{equation}\label{eq:rel-central-square}
 \begin{split}
 D\bigl(\tau(s+2uz)\bigr)
  &=(s-2uz)(s+2uz)\\
  &=s^2-4u^2z^2=s^2.
 \end{split}
\end{equation}
The correction has total degree $-1$, since $uz$ has degree zero.
In pair notation the primitive is $(s\tau,2\tau u)$.
Its differential is $(s^2,0)$: the two contributions to the second
component are $-2su$ and $2su$.
The first component is a primitive of $s^2$ in $E$.
The second cancels its commutation term in the differential of $C$.
Under $T$, this pair gives a bimodule null-homotopy for
$s^2\operatorname{id}_P$.
For $s$ itself, the matrix equations would require $A=1$ and $sB=2$,
which has no polynomial solution.

The square $u$ has two distinct mathematical roles in the calculation.
Its class in $H^{-2}(E)$ is nonzero.
Its product with $z$ represents the nonzero scalar $t^r/2$ in central
degree zero. The identity $t^r=2uz$ records the obstruction to extending
the relation $t^r=0$ in $H^*(E)$ to the relative derived center.

\section{The first two relations}
\label{sec:rel-examples}

\begin{example}[The relation $t=0$]\label{ex:rel-r1}
Let $r=1$. The free differential graded algebra $E$ has $d\tau=t$ and
$H^*(E)=k[u]$, with the parameter $t$ acting as zero.
The cohomology classes are $1,u,u^2,\ldots$ in degrees $0,-2,-4,\ldots$.
For each $j$, the equation $d(\tau u^j)=tu^j$ identifies a primitive of $tu^j$.

In the central complex, $D\tau=t-2uz$ and $D(\tau z)=tz$.
For every $j\geq0$, the differential from degree $-2j-1$ to $-2j$ is
\[
 \begin{pmatrix}t&0\\-2&t\end{pmatrix}.
\]
Consequently each nonpositive even degree has the $k$-basis
$[u^j],[tu^j]$, while degree two has basis $[z]$.
The relations are $t^2=0$, $tz=0$, and $uz=t/2$.
In particular, the scalar $t$ is nonzero in central degree zero, despite
vanishing on $H^*(E)$.

The equation $t=2uz$ eliminates $t$ from the cohomology presentation and
gives a graded $k$-algebra isomorphism
\[
 \operatorname{HH}_R^*(E)\cong k[u,z]/(z^2).
\]
The retained $R$-action sends $t$ to $2uz$ in this presentation.
Indeed, substitution in either direction respects both defining
relations, so the two polynomial maps are inverse.
A primitive of the central scalar $t^2$ is given by
\[
 D\bigl(t\tau+2\tau uz\bigr)=t^2.
\]
Its first term has differential $t^2-2tuz$.
Its second has differential $2tuz$, which removes the defect.
\end{example}

\begin{example}[The relation $t^2=0$]\label{ex:rel-r2}
Let $r=2$. The cohomology algebra of $E$ is $k[t,u]/(t^2)$.
Each nonpositive even degree has basis $[u^j],[tu^j]$ over $k$.
The equations $d\tau=t^2$ and $d(t\tau)=t^3$ exhibit $t^2$ and $t^3$
as boundaries in degree zero.

The central differential is $D\tau=t^2-2uz$.
Its matrix in each nonpositive even degree is
\[
 \begin{pmatrix}t^2&0\\-2&t^2\end{pmatrix}.
\]
Thus each such degree has the four independent classes
\[
 [u^j],\quad[tu^j],\quad[t^2u^j],\quad[t^3u^j].
\]
Degree two has the two classes $[z],[tz]$.

The product satisfies
\[
 uz=\tfrac12t^2,\qquad u(tz)=\tfrac12t^3,\qquad t^2z=0.
\]
These relations and $z^2=0$ determine the products together with the
polynomial products of $t$ and $u$.

The scalar $t^3$ is nonzero centrally, and $t^4$ is zero.
The primitive is
\[
 D\bigl(t^2\tau+2\tau uz\bigr)=t^4.
\]
The term $2\tau uz$ contributes $2t^2uz$, cancelling the defect from
$t^2\tau$. Its coefficient is the same as for $r=1$, because it comes
from the graded sign associated with the odd generator $\tau$.
\end{example}

\section{The square of a primitive}
\label{sec:rel-square-invariant}

The class $[\tau^2]$ uses the product of two primitives of $t^r$.
We will prove that this class is independent of the primitive and is
preserved by algebra quasi-isomorphisms that fix the scalar action.
It will distinguish $E$ from its cohomology algebra even when the
comparison uses a sequence of such maps in either direction.

A morphism of unital differential graded $R$-algebras is a
quasi-isomorphism when its underlying chain map is one.
The zero-differential algebra $H^*(A)$ has the induced product and
the scalar action $a\mapsto[a1]$.

\begin{definition}\label{def:rel-formality}
A differential graded $R$-algebra $A$ is formal over $R$ if it can be
joined to $H^*(A)$ by a finite sequence of quasi-isomorphisms of unital
differential graded $R$-algebras, each arrow allowed in either direction.
Such a sequence is called a quasi-isomorphism zigzag.
\end{definition}

To exclude such a zigzag, we seek an invariant that vanishes at one
endpoint and not the other. Its vanishing must be preserved by
quasi-isomorphisms in either direction.

\begin{lemma}\label{lem:rel-square-invariant}
Let $A$ be a unital differential graded $R$-algebra with $H^{-1}(A)=0$.
Suppose $a\in R$ satisfies $[a1]=0$ in $H^0(A)$.
For any $h\in A^{-1}$ with $dh=a1$, the element $h^2$ is a cycle.
The class
\[
 \sigma_a(A)=[h^2]\in H^{-2}(A)
\]
is independent of $h$.
Every morphism of such $R$-algebras preserves this class whenever both
algebras satisfy the stated hypotheses.
\end{lemma}

\begin{proof}
The hypothesis $[a1]=0$ provides a primitive $h$ of degree minus one.
Since $a$ is a scalar, the Leibniz rule gives
\[
 d(h^2)=(dh)h-h(dh)=ah-ha=0.
\]

If $h'$ is another primitive of $a1$, then $h'-h$ is a degree-minus-one cycle.
The vanishing of $H^{-1}(A)$ gives $x\in A^{-2}$ with $h'=h+dx$.
Expansion yields
\[
 h'^2-h^2=(dx)h+h(dx)+(dx)^2.
\]
Its explicit primitive is
\begin{equation}\label{eq:rel-square-independence}
 h'^2-h^2=d\bigl(xh-hx+x\,dx\bigr).
\end{equation}

In fact,
\[
 \begin{split}
 d(xh)&=(dx)h+xa,\\
 d(hx)&=ax-h(dx),\\
 d(x\,dx)&=(dx)^2.
 \end{split}
\]
The scalar terms cancel because $xa=ax$.
This proves independence of the chosen primitive.

For a unital $R$-algebra morphism $F:A\to B$, the element $F(h)$
satisfies $dF(h)=a1$ and $F(h^2)=F(h)^2$.
It follows that $H^{-2}(F)(\sigma_a(A))=\sigma_a(B)$.
\end{proof}

\begin{theorem}\label{thm:rel-nonformality}
For every $r\geq1$, the algebra
$E=k[t]\langle\tau\rangle$ with $d\tau=t^r$ is not formal over $R=k[t]$.
\end{theorem}

\begin{proof}
Proposition~\ref{prop:rel-state-cohomology} gives $H^{-1}(E)=0$.
Take the scalar $a=t^r$ and its primitive $h=\tau$.
The invariant is
\[
 \sigma_{t^r}(E)=[\tau^2]=[u]\ne0.
\]
In the zero-differential algebra $H^*(E)$, the scalar $t^r$ is already
the zero element. It therefore has the primitive $h=0$, giving
$\sigma_{t^r}(H^*(E))=0$.

Suppose a quasi-isomorphism zigzag joined these algebras.
Every intermediate algebra would have $H^{-1}=0$, because
quasi-isomorphisms induce cohomology isomorphisms.
The scalar $t^r$ would also be a boundary in every intermediate algebra.
For a forward arrow this follows from preserving the scalar and its zero
class. For a backward arrow it follows from injectivity of the induced
map on $H^0$ and preservation of the unit.

Thus Lemma~\ref{lem:rel-square-invariant} applies at every step.
Along each arrow the induced isomorphism on $H^{-2}$ sends the scalar
square invariant to the scalar square invariant.
In either direction an isomorphism preserves whether that class is zero.
The nonzero class at $E$ cannot become the zero class at $H^*(E)$.
This contradiction excludes every such zigzag.
\end{proof}

The specified $R$-action is essential to the statement: it identifies
the same scalar $t^r$ in every algebra and every arrow.
The central exponent and the scalar square invariant use this fixed
action in different ways. The former tests whether a scalar map admits
a compatible null-homotopy on the resolved diagonal.
The latter tests whether the square of a primitive can survive a
multiplicative comparison with a zero-differential algebra.
For $E$, both calculations use the class of $u=\tau^2$.

\section{Changing the coefficients}
\label{sec:rel-coefficient-change}

Adjoining a primitive uses a specified scalar ring.
Changing that ring gives two questions: whether the universal property
persists, and whether cohomology commutes with the change of coefficients.
Let $R'$ be a commutative $R$-algebra, with scalar map
$\alpha:R\to R'$, and put $a=\alpha(t)$.
The element $a$ may be zero or may multiply nonzero elements to zero.

Extending coefficients gives the differential graded $R'$-algebra
\[
 E_{R'}=R'\otimes_R E=R'\langle\tau\rangle,
 \qquad d\tau=a^r.
\]
Its basis is still $1,\tau,\tau^2,\ldots$.

For a differential graded $R'$-algebra $A$, a morphism $f:E\to A$
over $R$ extends uniquely by
\[
 R'\otimes_R E\longrightarrow A,\qquad b\otimes x\longmapsto b f(x).
\]
Scalar balancing makes the formula well-defined, and centrality of the
scalars makes it multiplicative. Its degree and differential conditions
follow directly from those of $f$.
Restricting to $1\otimes E$ reverses the construction.
Consequently $E_{R'}$ represents primitives $h$ satisfying $dh=a^r1$
over the new scalar ring, with no restriction on $\alpha$.

The cohomology is a different question. Define the annihilator
\[
 \operatorname{ann}_{R'}(a^r)=\{b\in R':a^rb=0\}.
\]
It measures the possible failure of multiplication by $a^r$ to be
injective.

\begin{proposition}\label{prop:rel-specialization}
For every $j\geq0$, coefficient extension gives
\[
 \begin{split}
 H^{-2j}(E_{R'})&=R'/(a^r)\,[u^j],\\
 H^{-2j-1}(E_{R'})&=\operatorname{ann}_{R'}(a^r)\,\tau u^j,
 \qquad u=\tau^2.
 \end{split}
\]
All positive-degree cohomology groups vanish.
The comparison
\[
 H^*(E)\otimes_R R'\longrightarrow H^*(E_{R'}),
 \qquad [x]\otimes b\longmapsto[b\otimes x],
\]
is an isomorphism in every degree exactly when
$\operatorname{ann}_{R'}(a^r)=0$.
\end{proposition}

\begin{proof}
The word differential remains
$d(\tau^{2j})=0$ and $d(\tau^{2j+1})=a^r\tau^{2j}$.
Every even word is a cycle, and its boundaries have coefficients in
$(a^r)$. An odd element $b\tau u^j$ is a cycle exactly when $a^rb=0$.
There are no odd boundaries because the differential vanishes on the
even part. This proves the displayed groups.

The comparison is well-defined: replacing a cycle $x$ by $x+dy$
changes its image by the boundary $d(b\otimes y)$, and the scalar
balancing relations hold in the target.
In every nonpositive even degree it is the map
\[
 (R/(t^r))\otimes_R R'\longrightarrow R'/(a^r),
 \qquad [c]\otimes b\longmapsto[\alpha(c)b].
\]
Its inverse sends $[b]$ to $[1]\otimes b$.
This is independent of the representative because
$[1]\otimes a^r b=[t^r]\otimes b=0$.
The two composites are identities by balancing.

The source has no odd degrees, while the target's odd groups are copies
of $\operatorname{ann}_{R'}(a^r)$.
Hence the comparison is an isomorphism precisely under the stated
annihilator condition.
\end{proof}

An $R$-module $R'$ is flat if tensoring with it preserves every short
exact sequence of $R$-modules. Short exact has the same injection,
kernel, and surjection meaning as for complexes degree by degree.
Flatness is sufficient for the condition above.
Indeed, applying coefficient extension to
\[
 0\longrightarrow R\xrightarrow{\ t^r\ }R
 \longrightarrow R/(t^r)\longrightarrow0
\]
keeps multiplication by $a^r$ injective.
The proposition identifies the exact condition needed for this particular
complex, even when flatness of the entire coefficient module is unknown.

For $R'=k=R/(t)$, the image $a$ is zero.
Every word then survives, so
\[
 H^*(E_{k})=k\langle\tau\rangle,
 \qquad H^*(E)\otimes_R k=k[u].
\]
The comparison sends $u$ to $\tau^2$ and misses every odd word.
The equation defining the primitive extends to $d\tau=0$.
The zero differential produces the additional odd cohomology classes.
